% Part: history
% Chapter: biographies 
% Section: alonzo-church

\documentclass[../../../include/open-logic-section]{subfiles}

\begin{document}

\olfileid{his}{bio}{chu}

\olsection{અલોન્ઝો ચર્ચ}

\olphoto{church-alonzo}{અલોન્ઝો ચર્ચ}

અલોન્ઝો ચર્ચનો જન્મ 14 જૂન 1903ના રોજ
વૉશિંગ્ટન, ડી.સી.માં થયો હતો. નાનપણમાં
હવાઈ બંદૂક સાથેની ઘટનાને કારણે તેમની
એક આંખની દૃષ્ટિ જતી રહી. તેમણે 1920માં
કનેક્ટિકટમાં યુનિવર્સિટી-પૂર્વ શાળાનો અભ્યાસ પૂરો
કર્યો અને એ જ વર્ષે પ્રિન્સ્ટનમાં
યુનિવર્સિટીનો અભ્યાસ શરૂ કર્યો. 1927માં
તેમણે ડૉક્ટરેટનો અભ્યાસ પૂરો કર્યો.
કેટલાંક વર્ષ વિદેશમાં વિતાવ્યા પછી ચર્ચ
પ્રિન્સ્ટન પરત ફર્યા. તેઓ અત્યંત વિનમ્ર
અને ચીવટવાળા તરીકે જાણીતા હતા. બોર્ડ
પરનું તેમનું લખાણ એકદમ સ્વચ્છ હોતું,
અને મહત્વના કાગળોને સાચવવા તેઓ તેને
પારદર્શક ગુંદર, ડ્યુકો સિમેન્ટ, વડે
કાળજીપૂર્વક ઢાંકતા. શૈક્ષણિક કાર્ય
સિવાય તેમને વિજ્ઞાનકથાની સામયિકો
વાંચવાનો શોખ હતો; લેખનમાં ભૂલ દેખાય
તો સંપાદકોને લખતાં પણ તેઓ અચકાતા નહીં.

ચર્ચની શૈક્ષણિક સિદ્ધિઓ મોટી હતી.
તેમના વિદ્યાર્થીઓ સ્ટીફન ક્લીની અને
બાર્કલી રોસેર સાથે તેમણે અસરકારક
સંગણનીયતાનો સિદ્ધાંત, લૅમ્બડા કલન,
વિકસાવ્યો; એલન ટ્યુરિંગે ટ્યુરિંગ
મશીન વિકસાવ્યું તેનાથી આ કાર્ય સ્વતંત્ર
હતું. સંગણનીયતાની આ બે વ્યાખ્યાઓ
સમકક્ષ છે અને આજે જેને
\emph{ચર્ચ--ટ્યુરિંગ માન્યતા} કહે છે
તે તરફ દોરી જાય છે: પ્રાકૃતિક સંખ્યાઓ
પરનું વિધેય અસરકારક રીતે સંગણનીય હોય
ત્યારે અને ત્યારે જ જ્યારે તેને
ટ્યુરિંગ મશીન (અથવા લૅમ્બડા કલન)
વડે ગણી શકાય. તેમણે આજે
\emph{ચર્ચનું પ્રમેય} કહેવાતું પરિણામ
પણ સાબિત કર્યું: પ્રથમ-ક્રમનાં સૂત્રોનું
પ્રામાણ્ય નક્કી કરવાની નિર્ણય સમસ્યા
ઉકેલી શકાતી નથી.

ચર્ચે વૃદ્ધાવસ્થા સુધી પોતાનું કાર્ય
ચાલુ રાખ્યું. 1967માં તેઓ પ્રિન્સ્ટન
છોડી યુસીએલએ ગયા, જ્યાં 1990માં
નિવૃત્ત થયા ત્યાં સુધી પ્રોફેસર રહ્યા.
1 ઑગસ્ટ 1995ના રોજ 92 વર્ષની ઉંમરે
તેમનું અવસાન થયું.

\begin{reading}
ચર્ચના ટૂંકા જીવનચરિત્ર માટે
\citet{EndertonND} જુઓ. લૅમ્બડા કલન
અને નિર્ણય સમસ્યા (ચર્ચની માન્યતા)
વિશે ચર્ચનાં મૂળ લખાણો
\citet{Church1936,Church1936a} છે.
\citet{Aspray1984}માં 1930ના દાયકાના
પ્રિન્સ્ટન ગણિતસમુદાય વિશે ચર્ચની
મુલાકાત નોંધાઈ છે. ચર્ચે 1936થી
1979 સુધી \emph{Journal of
Symbolic Logic}માં પુસ્તક-સમીક્ષાઓની
શ્રેણી લખી. તે બધી જૉન મૅકફાર્લેનની
વેબસાઇટ પર સંગ્રહિત છે
\citep{MacFarlane2015}.
\end{reading}
\end{document}
